\section{Downstream results}
\label{ch:downstream}

Potential extensions of the project once the main theorem is in place. They are listed
roughly in order of increasing formalization cost. The first two stay within
finite-dimensional linear algebra plus elementary functional analysis, while the
operator-algebraic consequences require von Neumann algebra infrastructure that
largely does not yet exist in Mathlib and would constitute a major sub-project of
independent value.

\subsection{Uncomputability of the value}

Corollary~\ref{cor:value-uncomputable} is already a headline consequence requiring
nothing further: no algorithm can approximate the (synchronous or quantum) value of a
nonlocal game to within an additive constant $c < \frac14$, since an estimate of additive
error $c$ separates value $1$ from value $\frac12$ exactly when $c < \frac14$ (at
$c = \frac14$ the estimate $\frac34$ is consistent with both). The threshold rises to the
trivial $\frac12$ once the gap is amplified: repeating the game drives the soundness value
below any $\delta > 0$ while preserving value $1$
(Theorem~\ref{thm:direct-repetition-q}), at a cost polynomial in the description, so for
every $c < \frac12$ no algorithm approximates the value to within additive error $c$. That
stronger form depends on the repetition theorem as well as on the main theorem.

\subsection{Semidecision from both sides}
\label{sec:ds-semidecision}

The non-explicit route mentioned below, and the explicit separating game, both rest on the
same pair of procedures: the entangled value is semidecidable from below, the commuting value
from above. Neither is hard, but both have hypotheses a paraphrase loses, so they are stated.

\effortMedium

\begin{lemma}[Lower semidecision of the entangled value]
\label{lem:val-lower-re}
\ledgernode{1.8.1}
\uses{def:q-value,lem:value-lower-approx}
For a nonlocal game $\game$ with rational data and rational $r$, the relation
$\valstar(\game) > r$ is recursively enumerable. Enumerate dimensions together with vectors
and POVMs whose entries are exact algebraic numbers; these are effectively enumerable and
dense among all finite-dimensional strategies, the payoff is continuous, so a strict
inequality is witnessed at some enumerated point, and exact algebraic comparison certifies it.
No supremum is assumed to be attained, and the supremum over $C_q$ agrees with the one over
$C_{qa}$ by continuity.
\end{lemma}

\effortHard

\begin{theorem}[NPA convergence and compactness]
\label{thm:npa-convergence}
\ledgernode{1.8.2, 1.8.2.6}
\uses{def:commuting-value,lem:npa-dilation,lem:npa-levels,lem:npa-moment-bound,lem:npa-diagonal,lem:npa-gns}
Fix a finite bipartite scenario. The normalized full-word level-$k$ NPA feasible sets have
uniformly bounded moments, their correlation projections $K_k$ are compact and decreasing, and
$\bigcap_k K_k = C_{qc}$, with the POVM and projective definitions of $C_{qc}$ equivalent. In
particular $C_{qc}$ is compact, every finite game payoff attains its maximum over it, and
$C_{qa} \subseteq C_{qc}$. No finite-dimensional attainment and no convergence rate is
asserted. The five lemmas below are its ingredients.
\end{theorem}

\begin{lemma}[Common projective dilation of commuting POVMs]
\label{lem:npa-dilation}
\ledgernode{1.8.2.1}
\uses{def:commuting-value}
Commuting finite-outcome POVMs admit a correlation-preserving common \emph{projective}
dilation. For $Vh = (\sqrt{A_a} h)_a$ the operator
$U = \begin{pmatrix} 0 & V^* \\ V & \Id - VV^* \end{pmatrix}$ is a self-adjoint
unitary on $\mH \oplus \mH^r$ whose conjugated coordinate projections compress to $A_a$, and
all questions share one embedding. Across the two players the block entries stay in commuting
algebras, so iterating the dilation preserves both commutation and the joint distribution. Zero
effects and arbitrary Hilbert space dimension are allowed.
\end{lemma}

\begin{lemma}[The full-word levels, and that they are effectively enumerable]
\label{lem:npa-levels}
\ledgernode{1.8.2.2}
\uses{def:commuting-value}
For a finite Bell scenario, level $k$ of the full-word NPA hierarchy consists of the normalized
positive semidefinite moment matrices on words of length at most $k$, subject to the
contextual projection, completeness and cross-commutation relations of degree at most $2k$,
with two-letter moments equal to the correlation coordinates. These constraints are finite and
effectively enumerable, every $C_{qc}$ correlation is feasible at every level, the levels are
nested, and this redundant presentation agrees with the standard reduced one by
degree-preserving linear substitutions.
\end{lemma}

\begin{lemma}[Moment entries are bounded, so the levels are compact]
\label{lem:npa-moment-bound}
\ledgernode{1.8.2.3}
\uses{lem:npa-levels}
Every feasible level-$k$ moment satisfies $\abs{L_k(w)} \leq 1$ for $\abs{w} \leq 2k$: the
identity $\sum_a \Gamma[e_a v, e_a v] = \Gamma[v,v]$ with positivity bounds each
leading-generator diagonal by its suffix diagonal, and induction together with the positive
semidefinite $2 \times 2$ minors bounds every entry. Hence each finite-level feasible set, and
its correlation projection, is compact.
\end{lemma}

\begin{lemma}[A diagonal subsequence, from certificates that need not be compatible]
\label{lem:npa-diagonal}
\ledgernode{1.8.2.4}
\uses{lem:npa-moment-bound}
If a correlation is feasible at every level, a countable diagonal subsequence of the bounded
certificates --- which need \emph{not} be compatible across levels, and this is the step where
a paraphrase overclaims --- defines a normalized positive functional on the polynomial
$*$-algebra modulo the measurement relations. Every polynomial positivity constraint and every
finite contextual relation survives the limit.
\end{lemma}

\begin{lemma}[GNS: a positive functional gives a commuting projective representation]
\label{lem:npa-gns}
\ledgernode{1.8.2.5}
\uses{lem:npa-diagonal}
A normalized positive functional on the measurement $*$-algebra yields a separable commuting
projective representation. Quotient by the zero-norm subspace and complete; left multiplication
by a projection generator $g$ is a contraction because
$\norm{u}^2 - \norm{gu}^2 = \norm{(\Id - g)u}^2$, so the null subspace is invariant, and the
PVM and cross-commutation identities extend from the dense polynomial quotient. The
distinguished unit vector is the class of $1$.
\end{lemma}

\effortEasy

\begin{lemma}[Decidability of the first-order theory of real closed fields]
\label{lem:rcf-decision}
\ledgernode{1.8.3}
A first-order sentence over real closed fields with rational coefficients is decidable; in
particular so is feasibility of a finite conjunction of polynomial equalities and weak or
strict inequalities over $\R$.
\end{lemma}

\paragraph{Source.} Tarski; the effective form used here is Basu, Pollack and
Roy~\cite{BPR06}. This is an import: the companion ledger records it as \emph{admitted}, one
of only three assumptions in the campaign (Remark~\ref{rem:admitted-nodes}), and it is the
only one that this chapter, rather than the main theorem, is responsible for.

\effortMedium

\begin{lemma}[Upper semidecision of the commuting value]
\label{lem:valco-upper-re}
\ledgernode{1.8.4, 1.8.4.1}
\uses{thm:npa-convergence,lem:rcf-decision,def:commuting-value}
For a game $\game$ with rational data and rational $r$, the relation $\valco(\game) < r$ is
recursively enumerable: at each full-word NPA level decide exact rational semialgebraic
feasibility of payoff $\geq r$ using Lemma~\ref{lem:rcf-decision}. Some level is infeasible
exactly when $\valco(\game) < r$; otherwise Theorem~\ref{thm:npa-convergence} assembles a
commuting-operator strategy of payoff $\geq r$. Dovetailing over rational $r < 1$ halts
exactly when $\valco(\game) < 1$. No efficiency claim and no numerical SDP certificate is
involved --- the decision is exact, and succinct descriptions are expanded first.
\end{lemma}

\subsection{Separation of quantum and commuting values}

\begin{theorem}[Separation]
\label{thm:separation}
\ledgernode{1.8.8, 1.8.9, 1.8.10, 1.8.11}
\uses{thm:halting,def:commuting-value,def:q-value,lem:valco-upper-re,thm:compression}
There is an explicit synchronous game $\game^{\sep}$ with
$\valstar(\game^{\sep}) \leq \frac12$ and $\valco(\game^{\sep}) = 1$.
\end{theorem}

Sketch: the bound on $\valstar$ is the soundness clause of Theorem~\ref{thm:halting}
verbatim, now that the clause is stated in $\valstar$; it does not follow from a $\synval$
bound. Apply that reduction to a machine that searches
for a proof of its own non-halting (or any non-halting machine whose game retains a
perfect commuting-operator strategy), and track through the pipeline that
completeness holds not just for PCC strategies but for
commuting-operator (equivalently, tracial) strategies. This \emph{commuting
completeness} property is an additional obligation on
Sections~\ref{sec:ps-introspection}--\ref{sec:ps-compression} that the blueprint
should make explicit when those sections are worked out; it is the only new
ingredient beyond the main theorem. An alternative route: if $\valstar =
\valco$ for all games, then combining Lemma~\ref{lem:value-lower-approx} (value
approximable from below) with the NPA hierarchy~\cite{NPA08} (commuting value
approximable from above) would make the value computable, contradicting
Corollary~\ref{cor:value-uncomputable}; this gives a non-explicit separation with
much less work, and is the recommended first target.

\subsection{Failure of Tsirelson's problem}

\begin{corollary}[Tsirelson's problem, negative answer]
\label{cor:tsirelson}
\ledgernode{1.8.5}
\uses{thm:separation,lem:val-lower-re,lem:valco-upper-re,cor:value-uncomputable}
There exist finite question and answer sets for which the closure $C_{qa}$ of the set
of finite-dimensional quantum correlations is strictly contained in the set $C_{qc}$
of commuting-operator correlations.
\end{corollary}

Immediate from Theorem~\ref{thm:separation}: the value gap is witnessed by a
correlation in $C_{qc}$ at distance bounded away from $C_{qa}$. Requires only the
definitions of the correlation sets. (Non-closure of the unclosed set $C_q$ itself is
Slofstra's earlier theorem~\cite{Slo19}, which also follows, with quantitative bounds,
from the machinery here.)

\subsection{Refutation of Connes' Embedding Problem}

\effortHard

\begin{lemma}[Kirchberg's forward implication]
\label{lem:kirchberg-forward}
\ledgernode{1.8.6.1}
If CEP has an affirmative answer then the minimal and maximal C$^*$-tensor norms coincide on
$C^*(F_\infty) \ot_{\mathrm{alg}} C^*(F_\infty)$, where $F_\infty$ is the free group on
countably many generators and $C^*$ denotes the full unital group algebra.
\end{lemma}

\paragraph{Source.} Kirchberg; see~\cite{Oza13}. This is the second of the two imports this
chapter carries: the companion ledger records it as \emph{admitted} (node $1.8.6.1$), and it is
the only step of the chain below that is not proved there.

\begin{lemma}[Finitely many generators]
\label{lem:cep-free-group}
\ledgernode{1.8.6.2}
\uses{lem:kirchberg-forward}
Conditional on CEP, $\min = \max$ on
$C^*(F_k) \ot_{\mathrm{alg}} C^*(F_k)$ for every finite $k \geq 1$.
\end{lemma}

\begin{lemma}[A completely positive section onto the Bell algebra]
\label{lem:bell-algebra-section}
\ledgernode{1.8.6.3}
For $A = \ast_{x=1}^{k} \C^m$ and $B = C^*(F_k)$ with $k \geq 1$ and $m \geq 2$, the
canonical quotient $q : B \to A$ has a unital completely positive section $\Psi$, obtained as
the $(0,0)$ corner of an explicitly constructed unital $*$-homomorphism
$\gamma : A \to M_m(B)$. In particular $q \Psi = \mathrm{id}$ on all of $A$.
\end{lemma}

\begin{lemma}[Transport to the Bell algebra]
\label{lem:cep-bell-algebra}
\ledgernode{1.8.6.4}
\uses{lem:cep-free-group,lem:bell-algebra-section}
Conditional on CEP, the minimal and maximal tensor norms coincide on
$A \ot_{\mathrm{alg}} A$ for every finite Bell algebra $A = \ast_{x=1}^{k} \C^m$,
$k \geq 1$, $m \geq 2$.
\end{lemma}

\begin{lemma}[Commuting correlations restrict states on the maximal tensor product]
\label{lem:commuting-restricts-state}
\ledgernode{1.8.6.5}
\uses{def:commuting-value,lem:cep-bell-algebra}
Every commuting-operator correlation in a finite bipartite scenario with common parameters
$(k,m)$ is the restriction to the coordinate projections of a state on
$A \ot_{\max} A$. Under CEP that same functional is a state on $A \ot_{\min} A$.
\end{lemma}

\begin{lemma}[Vector states are dense on finitely many coordinates]
\label{lem:spatial-coordinates}
\ledgernode{1.8.6.6}
On any finite collection of self-adjoint coordinate observables in a faithfully represented
spatial C$^*$-tensor product, every state is approximable by finite convex combinations of
unit-vector states in that same spatial representation.
\end{lemma}

\begin{lemma}[States on the spatial product give $C_{qa}$ correlations]
\label{lem:spatial-values-cqa}
\ledgernode{1.8.6.7}
\uses{lem:spatial-coordinates}
The coordinate values of any state on a spatial tensor product of finite Bell algebras belong
to $C_{qa}$, the closure of the finite-dimensional tensor-product POVM correlations.
\end{lemma}

\begin{lemma}[CEP forces equality of the correlation sets]
\label{lem:cep-implies-equal}
\ledgernode{1.8.6, 1.8.6.8}
\uses{def:commuting-value,lem:commuting-restricts-state,lem:spatial-values-cqa}
If every separable II$_1$ factor embeds into an ultrapower of the hyperfinite II$_1$ factor,
then $C_{qa} = C_{qc}$ in every finite bipartite scenario, including unequal alphabet sizes.
Consequently exhibiting one finite-scenario separation of $C_{qa}$ from $C_{qc}$ refutes CEP.
\end{lemma}

\begin{corollary}[Connes' Embedding Problem, negative answer]
\label{cor:cep}
\ledgernode{1.8.7}
\uses{cor:tsirelson,lem:cep-implies-equal}
There is a separable II$_1$ factor (indeed, a tracial von Neumann algebra generated by
finitely many projections) that does not embed into an ultrapower
$\cal{R}^{\omega}$ of the hyperfinite II$_1$ factor.
\end{corollary}

\paragraph{Source for the implication.} Tsirelson's problem is equivalent to CEP by
Fritz~\cite{Fri12} and Junge et al.~\cite{JNPPSW11}, with the converse direction
completed by Ozawa~\cite{Oza13}. In the synchronous framework the implication is
particularly direct: a synchronous correlation is precisely given by projections in a
tracial von Neumann algebra summing to $1$~\cite{PSSTW16,KPS18}, and $C_{qa}$-type
correlations are exactly those arising from $\cal{R}^\omega$-embeddable algebras, so
Corollary~\ref{cor:tsirelson} for synchronous correlations yields a non-embeddable
algebra.

\paragraph{Comments.} This is the deepest formalization challenge on the list: it
needs tracial von Neumann algebras, ultraproducts, and the synchronous-correlation
correspondence. The synchronous framework adopted in this blueprint
(Section~\ref{sec:departures}) was chosen partly because it makes this bridge as
short as it can be. A reasonable staging: (i) define $C_{qa}, C_{qc}$ and synchronous
correlations; (ii) prove the synchronous correlation--tracial algebra correspondence;
(iii) formalize the ultraproduct $\cal{R}^\omega$ and the equivalence. Step (iii)
alone is a landmark for operator algebras in Lean.

\subsection{Failure of Kirchberg's QWEP conjecture}

\begin{corollary}[QWEP]
\label{cor:qwep}
\uses{cor:cep}
Kirchberg's QWEP conjecture fails: there is a C$^*$-algebra that is not a quotient of
a C$^*$-algebra with the weak expectation property; equivalently,
$C^*(F_2) \ot_{\min} C^*(F_2) \neq C^*(F_2) \ot_{\max} C^*(F_2)$.
\end{corollary}

Equivalent to CEP by Kirchberg's theorem (see~\cite{Oza13}); no new input from the
project, but the equivalence itself is another substantial operator-algebra
formalization.

\subsection{Further directions}

\begin{remark}
\label{rem:further}
Other candidate consequences to record as the project matures: undecidability of the
existence of \emph{perfect} strategies for synchronous games and its group-theoretic
reformulations (via the synchronous algebra of a game); quantitative non-closure of
$C_q$ and rates for the NPA hierarchy; uncomputability results for entanglement
requirements ($\Ent(\cdot, \frac12)$ grows faster than any computable function along
the halting games), which need the entanglement form of the pipeline
(Remark~\ref{rem:entanglement-form}) rather than the value form adopted here; the commuting-operator
counterpart $\MIPco = \coRE$~\cite{Lin23,Lin25}, for which direct parallel repetition
for commuting-operator strategies (Theorem~\ref{thm:direct-repetition-co}) and Lin's
tracial density theorem (Theorem~\ref{thm:tracial-density}) are already formalized; and,
more speculatively, connections
to the undecidability of the spectral gap and to the existence of non-hyperlinear
groups, the latter being \emph{not} implied by CEP's failure and best recorded as an
open problem.
\end{remark}

\subsection{Provenance of this chapter}

\begin{remark}[Stage 1.8 of the ledger, and how much of it has been challenged]
\label{rem:downstream-nodes}
\ledgernode{1.8}
\uses{thm:separation,cor:tsirelson,cor:cep,lem:cep-implies-equal,lem:valco-upper-re,lem:rcf-decision,rem:admitted-nodes}
The statements of this chapter were roadmap sketches until September 2026, when the companion
campaign opened a stage for them: stage $1.8$, twenty-seven nodes, which the statements above
now account for one for one. Three things about that stage are worth carrying here rather than
leaving in the other repository.

First, it is where two of the campaign's three remaining admissions live ---
Lemma~\ref{lem:rcf-decision} (node $1.8.3$) and the forward Kirchberg implication inside
Lemma~\ref{lem:cep-implies-equal} (node $1.8.6.1$). Neither is reached by the main theorem:
they are assumptions of the \emph{corollaries}, which is the right place for them, and
Remark~\ref{rem:admitted-nodes} keeps the main theorem's own surface to a single classical
algorithmic result.

Second, the separating game is more delicate than Theorem~\ref{thm:separation}'s statement
suggests, and the nodes say where. The construction needs effective bounded self-reference at a
fixed $C_0 \geq 2$ through one canonical specialization-and-timeout compiler (node $1.8.8$);
the commuting value is $1$ because if the upper semidecider ever halted, an early-accept branch
would supply a genuine one-dimensional perfect commuting strategy (node $1.8.9$); and
$\valstar \leq \frac12$ is an entanglement-cost divergence argument along $n_{j+1} = 2^{n_j}$
--- \emph{not} an unattained-supremum argument (node $1.8.10$). Only the last of these is
visible in the sketch above, and a formalization will meet all three.

Third, and this is the part to be honest about: these nodes have had far less adversarial
scrutiny than the rest of the ledger. Two of the twenty-seven carry a challenge, against $85\%$
of the nodes in the original campaign state, and the adversarial round of 2026-09-14 --- which
re-examined seventy-five nodes and found one real defect
(Remark~\ref{rem:sync-transfer-marginal}) --- worked entirely within stage $1.2$ and did not
touch stage $1.8$ at all. Node $1.8$'s own statement says its derivations ``remain pending
adversarial review''. So the annotations in this chapter record that a decomposition exists and
that it has been checked once, not that it has been attacked; treat the statements here as
better-specified sketches rather than as verified mathematics, and expect the numbered details
above to move.
\end{remark}
